\chapter{How to use this site}
\label{ch:intro}

This site lets a mathematician audit, \emph{without reading any proof
or learning any proof assistant}, what exactly has been
machine-checked in the
\href{https://github.com/LLM4Rocq/digraph-theory}{\texttt{digraph-theory}}
library:

\begin{itemize}
\item \textbf{for every $k \in \{3,4,5\}$ there are infinitely many
  $k$-$\bar\omega$-critical tournaments} (Conjecture 5.10 of
  Aboulker--Aubian--Charbit--Lopes, arXiv:2310.04265, at every value
  of $k$ for which it is humanly known) ---
  \ref{res:conjecture_5_10_at_345};
\item \textbf{Question 5.9 of the same paper fails at each such $k$}
  --- \ref{res:question_5_9_fails_at_k3},
  \ref{res:question_5_9_fails_at_k4},
  \ref{res:question_5_9_fails_at_k5};
\item \textbf{the $\delta = 3$ case of the Cheng--Keevash path
  conjecture}: every oriented digraph with minimum out-degree $\ge 3$
  contains a directed path of length 6 ---
  \ref{res:ck_conj1_delta3};
\item the general tournament theory behind them (substitution bound,
  vertex-transitive criticality, domination,
  uniqueness of the 2-critical tournament).
\end{itemize}

\paragraph{The division of labour.} The Rocq kernel has checked every
proof; each result here reports \emph{axiom-free} status
(\code{Print Assumptions}: closed under the global context). Proofs
are therefore not shown. What a human must check is that the
\emph{statements} say what the informal mathematics means --- so every
result page shows the verbatim formal statement, a plain-mathematics
decoding, and machine-generated \code{uses} links to \emph{every}
definition the statement depends on (the Dictionary,
\ref{ch:dictionary}). Use the dependency graph to explore: a result's
ancestors are exactly the definitions you must read to trust it.

\paragraph{Generated, not curated.} The dependency edges are computed
from the compiled library by \code{scripts/statement\_closure.py};
continuous integration fails if a statement's closure ever contains a
definition without a Dictionary entry, or if a quoted text drifts
from the source.

\paragraph{Conventions.} Statements are quoted verbatim. Boolean
connectives (\code{==}, \code{\&\&}, \code{[forall ...]}) are
MathComp's decidable counterparts of the usual logical ones ---
equivalent to them on the finite structures used here. The circulant
families use the index convention $m =$ \code{m'.+1},
$n =$ \code{m.*2.+1} with hypothesis \code{3 <= m'.+1}: $n$ ranges
exactly over the odd numbers $\ge 7$. A statement quoted from inside
a \code{Section} is implicitly universally quantified over the
section's parameters. A companion PDF with the same content plus a
notation primer and the full declaration catalog is available as
\href{https://llm4rocq.github.io/digraph-theory/digraph_formal.pdf}{\texttt{digraph\_formal.pdf}};
the per-file API documentation (coqdoc) is under
\href{https://llm4rocq.github.io/digraph-theory/doc}{\texttt{/doc}}.
